%% abtex2-modelo-projeto-pesquisa.tex, v<VERSION> laurocesar
%% Copyright 2012-<COPYRIGHT_YEAR> by abnTeX2 group at http://www.abntex.net.br/ 
%%
%% This work may be distributed and/or modified under the
%% conditions of the LaTeX Project Public License, either version 1.3
%% of this license or (at your option) any later version.
%% The latest version of this license is in
%%   http://www.latex-project.org/lppl.txt
%% and version 1.3 or later is part of all distributions of LaTeX
%% version 2005/12/01 or later.
%%
%% This work has the LPPL maintenance status `maintained'.
%% 
%% The Current Maintainer of this work is the abnTeX2 team, led
%% by Lauro César Araujo. Further information are available on 
%% http://www.abntex.net.br/
%%
%% This work consists of the files abntex2-modelo-projeto-pesquisa.tex
%% and abntex2-modelo-references.bib
%%

% ------------------------------------------------------------------------
% ------------------------------------------------------------------------
% abnTeX2: Modelo de Projeto de pesquisa em conformidade com 
% ABNT NBR 15287:2011 Informação e documentação - Projeto de pesquisa -
% Apresentação 
% ------------------------------------------------------------------------ 
% ------------------------------------------------------------------------

\documentclass[
	% -- opções da classe memoir --
	%article,
	12pt,				% tamanho da fonte
	openright,			% capítulos começam em pág ímpar (insere página vazia caso preciso)
	twoside,			% para impressão em recto e verso. Oposto a oneside
	a4paper,			% tamanho do papel. 
	% -- opções da classe abntex2 --
	%chapter=TITLE,		% títulos de capítulos convertidos em letras maiúsculas
	%section=TITLE,		% títulos de seções convertidos em letras maiúsculas
	%subsection=TITLE,	% títulos de subseções convertidos em letras maiúsculas
	%subsubsection=TITLE,% títulos de subsubseções convertidos em letras maiúsculas
	% -- opções do pacote babel --
	english,			% idioma adicional para hifenização
	brazil,				% o último idioma é o principal do documento
	]{abntex2}

% ---
% PACOTES
% ---

% ---
% Pacotes fundamentais 
% ---
\usepackage{lmodern}			% Usa a fonte Latin Modern
\usepackage[T1]{fontenc}		% Selecao de codigos de fonte.
\usepackage[utf8]{inputenc}		% Codificacao do documento (conversão automática dos acentos)
\usepackage{indentfirst}		% Indenta o primeiro parágrafo de cada seção.
\usepackage{color}				% Controle das cores
\usepackage{graphicx}			% Inclusão de gráficos
\usepackage{microtype} 			% para melhorias de justificação
\usepackage{amsmath}
% ---

% ---
% Pacotes adicionais, usados apenas no âmbito do Modelo Canônico do abnteX2
% ---
\usepackage{lipsum}				% para geração de dummy text
% ---

% ---
% Pacotes de citações
% ---
\usepackage[brazilian,hyperpageref]{backref}	 % Paginas com as citações na bibl
\usepackage[alf]{abntex2cite}	% Citações padrão ABNT

% --- 
% CONFIGURAÇÕES DE PACOTES
% --- 

% ---
% Configurações do pacote backref
% Usado sem a opção hyperpageref de backref
\renewcommand{\backrefpagesname}{Citado na(s) página(s):~}
% Texto padrão antes do número das páginas
\renewcommand{\backref}{}
% Define os textos da citação
\renewcommand*{\backrefalt}[4]{
	\ifcase #1 %
		Nenhuma citação no texto.%
	\or
		Citado na página #2.%
	\else
		Citado #1 vezes nas páginas #2.%
	\fi}%
% ---
% Macro para nome da universidade
\newcommand{\unb}{Universidade de Brasília - UNB}
\newcommand{\cic}{Departamento de Ciência da Computação - CIC}

% ---
% Informações de dados para CAPA e FOLHA DE ROSTO
% ---
\titulo{\textbf{Relatório Final de Lógica de Computação 1}}
\autor{Wagner Clemente Coelho Batalha}
\local{Brasília - Brasil}
\data{\today}
\instituicao{%
  \unb
  \par
  \cic}
\tipotrabalho{Relatório final}

\renewcommand{\imprimircapa}{%
	\begin{capa}%
		\center
		\ABNTEXchapterfont\Large \cic \par
		\vspace*{1cm} 
		{\ABNTEXchapterfont\large\imprimirautor}
		\vfill
		\begin{center}
		\ABNTEXchapterfont\bfseries\LARGE\imprimirtitulo
		\end{center}
		\vfill
		\large\imprimirlocal \par
		\large\imprimirdata
		\vspace*{1cm}
		\begin{figure}[t!] \center \includegraphics[width=0.75\linewidth]{img/as_comp_PB.jpg} \end{figure}
	\end{capa}
}
% O preambulo deve conter o tipo do trabalho, o objetivo, 
% o nome da instituição e a área de concentração 
\preambulo{Relatório final de Lógica de Computação 1}
% ---

% ---
% Configurações de aparência do PDF final

% alterando o aspecto da cor azul
\definecolor{blue}{RGB}{41,5,195}

% informações do PDF
\makeatletter
\hypersetup{
     	%pagebackref=true,
		pdftitle={\@title}, 
		pdfauthor={\@author},
    	pdfsubject={\imprimirpreambulo},
	    pdfcreator={LaTeX with abnTeX2},
		pdfkeywords={tdah}{educacao}{computacao}{projeto de pesquisa}, 
		colorlinks=true,       		% false: boxed links; true: colored links
    	linkcolor=blue,          	% color of internal links
    	citecolor=blue,        		% color of links to bibliography
    	filecolor=magenta,      		% color of file links
		urlcolor=blue,
		bookmarksdepth=4
}
\makeatother
% --- 

% --- 
% Espaçamentos entre linhas e parágrafos 
% --- 

% O tamanho do parágrafo é dado por:
\setlength{\parindent}{1.3cm}

% Controle do espaçamento entre um parágrafo e outro:
\setlength{\parskip}{0.2cm}  % tente também \onelineskip

% ---
% compila o indice
% ---
\makeindex
% ---

% ----
% Início do documento
% ----
\begin{document}

% Seleciona o idioma do documento (conforme pacotes do babel)
%\selectlanguage{english}
\selectlanguage{brazil}

% Retira espaço extra obsoleto entre as frases.
\frenchspacing 

% ----------------------------------------------------------
% ELEMENTOS PRÉ-TEXTUAIS
% ----------------------------------------------------------
% \pretextual

% ---
% Capa
% ---
\imprimircapa
% ---

% ---
% Folha de rosto
% ---
% \imprimirfolhaderosto
% ---

% ---
% NOTA DA ABNT NBR 15287:2011, p. 4:
%  ``Se exigido pela entidade, apresentar os dados curriculares do autor em
%     folha ou página distinta após a folha de rosto.''
% ---

% ---
% inserir o sumario
% ---
% \pdfbookmark[0]{\contentsname}{toc}
% \tableofcontents*
% \cleardoublepage
% ---


% ----------------------------------------------------------
% ELEMENTOS TEXTUAIS
% ----------------------------------------------------------
\textual

% Eu estou escrevendo os textos separando cada parte em um arquivo tex própio e
% depois inserindo neste documento com o \input{} para facilidade de
% leitura, escrita do texto e organização em geral.
\chapter{Introdução e contextualização do problema}

Algoritmo é o conceito mais importante na ciência da computação. A criação deles é essencial para todo tipo de solução de problema, automação, análise e captura de dados, comunicação e quase tudo relacionado com computação.

Entretanto a criação de algoritmos ficam a cargo de um ou mais programadores que podem introduzir erros, não detectar casos periféricos, ou até mesmo durante a fase de execução de código descobrir que a própria linguagem de programação utilizada introduz erros que não poderiam ter sido previstos.

Por isso que existem testes de unidade e integração. Ambos servem para detectar erros no código fonte que foram introduzidos pelo programador. Infelizmente o teste de unidade depende da criação de casos de teste pelo próprio programador. Apesar disto não impedir o aumento da confiabilidade do código, testes de unidade são limitados á linguagem de programação que foi utilizado e da capacidade do programador de adiar possíveis casos onde o algoritmo escrito pode falhar. Outro problema se revela quando é utilizado funções, métodos ou API de outras fontes, tornando necessário que se confie na ausência de falhas e dos testes dos algoritmos escritos por terceiros.

Portanto há a necessidade de verificação de algoritmos de forma universal, não apenas em uma só linguagem ou repositório de código, mas sim no contexto matemático de verificação, explorando lógica dedutiva para chegar na formalização de algoritmos sem ter que analisar exaustivamente de forma limitada ou finita (como no caso de testes de unidade) o código.

Portanto há diversos provadores lógicos para verificação, como o PVS (\textit{Prototype Verification System}) que é usado neste projeto.

\section{Formalização do algoritmo de ordenação Heapsort}

Neste trabalho será provado, através do PVS, o correto funcionamento do Heapsort. Com esta prova podemos implicar que se for escrito o algoritmo em qualquer linguagem de programação (desde que seja idêntico ao utilizado neste projeto), ele irá funcionar com garantia lógica.

\section{Explicação das técnicas de solução}

É importante notar que o PVS implementa lógica de sequente de Gentzen dentro de um ambiente de prova iterativo de forma indireta através de comandos de prova, notando que muitos desses comandos  executam de forma automática as regras do sequente.



\chapter{Especificação do problema e explicação do método de solução}

Há dois arquivos na prova do heapsort: o \textit{heapsort.pvs} e o \textit{sorting\_seq.pvs}.

O arquivo \textit{heapsort.pvs} contêm o algoritmo de ordenação divido em várias partes. Cada parte é provada separadamente e apenas no final do arquivo é provado o algoritmo de forma universal.

O arquivo \textit{sorting\_seq.pvs} tem diversas funções para verificação de estados de lista como encontrar a quantidade de ocorrências de um item em uma lista, descobrir se uma lista é uma permutação de outra ou se a lista está ordenada. Este arquivo também tem vários teoremas e corolários utilizados nas provas do arquivo \textit{heapsort.pvs}.

\section{Provando heapsort}

O teorema \textit{heapsort\_works} é definido da seguinte forma:

\begin{verbatim}
heapsort\_works : THEOREM 
FORALL (h) : is\_sorted?(heapsort(h)) AND permutations(h, heapsort(h))
\end{verbatim}

A função acima descreve que para toda lista finita h (O tipo de h é definido no começo do arquivo pvs, não necessitando a explicitação do tipo no teorema) heapsort retorna uma lista ordenada e que a lista que ele retorna é uma permutação de h, sem perder dados.

A função \textit{permutations(h1, h2)} (onde h1 e h2 são do tipo lista finita) é definida em \textit{sorting\_seq.pvs} e retorna TRUE caso a lista h1 seja uma permutação de h2. 

Podemos deduzir que tal essa função tem algumas propriedades:

\begin{enumerate}
    \item \label{lemma2} Caso \textit{permutations(h1, h2)} seja verdadeiro, o tamanho de h1 deve ser o mesmo que o tamanho de h2, já que ambas listas são essencialmente as mesmas apenas com os lugares trocados;
    \item \label{reflex} Reflexividade. \textit{permutations(h1, h2)} deve retornar o mesmo valor caso o parâmetro passado seja trocado de lugar: \textit{permutations(h2, h1)};
    \item \label{symetry} Propriedade simétrica: \textit{permutations(h1, h1)} será sempre verdadeiro para qualquer h1;
    \item \label{trans} Propriedade transitiva: se \textit{permutations(h1, h2)} for verdadeiro e \textit{permutations(h2, h3)} também, \textit{permutations(h1, h3)} será verdadeiro.
\end{enumerate}

Todas as propriedades acima foram provadas em Corolários e Lemas. A propriedade \ref{lemma2} pode ser provado utilizando tanto uma prova sequencial como uma prova indutiva no comprimento da lista. 

As propriedades de reflexividade, simetria e transitividade foram concatenadas em apenas uma prova: \textit{permutations\_equiv}.

A função is\_sorted? por outro lado é definida e provada na biblioteca NASALib, que também serve de fonte de lemas, teoremas e corolários.

A prova de \textit{heapsort\_works} utiliza vários lemmas de ambas fontes e principalmente de provas das sub-funções do heapsort. Todas estas provas de propriedades facilitam a formalização de heapsort quando necessário através do corte e aplicação de regras proposicionais.


\section{Descrição da formalização}

% ----------------------------------------------------------
% ELEMENTOS PÓS-TEXTUAIS
% ----------------------------------------------------------
\postextual

% ----------------------------------------------------------
% Referências bibliográficas
% ----------------------------------------------------------
\bibliography{bibliografia}

% ----------------------------------------------------------
% Glossário
% ----------------------------------------------------------
%
% Consulte o manual da classe abntex2 para orientações sobre o glossário.
%
%\glossary

% ----------------------------------------------------------
% Apêndices
% ----------------------------------------------------------

% ---
% Inicia os apêndices
% ---
\begin{apendicesenv}

% Imprime uma página indicando o início dos apêndices
\partapendices

\lipsum[55-57]

\end{apendicesenv}
% ---


% ----------------------------------------------------------
% Anexos
% ----------------------------------------------------------

% ---
% Inicia os anexos
% ---
\begin{anexosenv}

% Imprime uma página indicando o início dos anexos
\partanexos

\end{anexosenv}

%---------------------------------------------------------------------
% INDICE REMISSIVO
%---------------------------------------------------------------------

\phantompart

\printindex


\end{document}

